% Figure: the wall-mounted jib crane of the case study. The assembled
% scene on the left; the chosen-body decomposition (boom and tie rod) on
% the right. Numbers match the worked case study in the text.
\begin{figure}[htbp]
\centering
\begin{tikzpicture}[
  member/.style={draw=engblack, line width=1.6pt},
  tie/.style={draw=engblue, line width=1.4pt},
  load/.style={draw=engred, -{Stealth}, very thick},
  scale=1.0,
]
% ===== Left: the assembled crane =====
\begin{scope}[xshift=0cm]
  % wall
  \draw[draw=engblack, very thick] (0,0) -- (0,4.4);
  \foreach \y in {0.3,0.7,...,4.2} {\draw[draw=enggray, thin] (0,\y) -- (-0.28,\y-0.2);}
  % pin at A (boom root)
  \coordinate (A) at (0,1.0);
  \coordinate (B) at (4.0,1.0);
  \coordinate (C) at (0,3.4);
  % boom A==B (horizontal)
  \draw[member] (A) -- (B);
  % tie rod C==B
  \draw[tie] (C) -- (B);
  % pin markers
  \fill[engblack] (A) circle (0.08);
  \draw[draw=engblack, fill=white] (A) circle (0.13);
  \fill[engblack] (C) circle (0.08);
  \draw[draw=engblack, fill=white] (C) circle (0.13);
  \fill[engblack] (B) circle (0.08);
  \draw[draw=engblack, fill=white] (B) circle (0.13);
  % node labels
  \node[font=\scriptsize, anchor=east] at (A) {$A$};
  \node[font=\scriptsize, anchor=east] at (C) {$C$};
  \node[font=\scriptsize, anchor=south west] at (B) {$B$};
  % hoisted load at the boom tip
  \draw[load] (B) -- ++(0,-1.4) node[anchor=north, font=\small, engred] {$W$};
  % tie-rod angle marker at B
  \draw[draw=enggray, thin] (3.4,1.0) arc (180:149:0.6);
  \node[font=\scriptsize] at (3.25,1.22) {$\beta$};
  % dimension hints
  \node[font=\scriptsize, anchor=north] at (2.0,1.0) {boom $L$};
  \node[font=\scriptsize, anchor=east, rotate=31] at (2.0,2.5) {tie rod};
  \node[font=\small, anchor=north] at (2.0,-0.1) {(a) the assembly};
\end{scope}
% ===== Right: FBD of the boom AB as a single body =====
\begin{scope}[xshift=7.2cm]
  \coordinate (A2) at (0,1.0);
  \coordinate (B2) at (4.0,1.0);
  \draw[member] (A2) -- (B2);
  \fill[engblack] (A2) circle (0.06);
  \fill[engblack] (B2) circle (0.06);
  \node[font=\scriptsize, anchor=north east] at (A2) {$A$};
  \node[font=\scriptsize, anchor=north] at (B2) {$B$};
  % pin reaction components at A
  \draw[draw=engred, -{Stealth}, very thick] (A2) -- ++(1.5,0)
    node[anchor=south, font=\scriptsize, engred] {$A_x$};
  \draw[draw=engred, -{Stealth}, very thick] (A2) -- ++(0,1.5)
    node[anchor=east, font=\scriptsize, engred] {$A_y$};
  % tie-rod force at B: directed up and to the left along the rod (tension)
  \draw[draw=engblue, -{Stealth}, very thick] (B2) -- ++(-2.0,1.2)
    node[anchor=south, font=\scriptsize, engblue] {$T$};
  \draw[draw=enggray, thin] (3.4,1.0) arc (180:149:0.6);
  \node[font=\scriptsize] at (3.22,1.22) {$\beta$};
  % hoisted load down at B
  \draw[load] (B2) -- ++(0,-1.4) node[anchor=north, font=\small, engred] {$W$};
  \node[font=\small, anchor=north] at (2.0,-0.1) {(b) FBD of the boom};
\end{scope}
\end{tikzpicture}
\caption[A wall-mounted jib crane and the free-body diagram of its
boom]{The wall-mounted jib crane of the case study. (a) The assembly: a
horizontal boom $AB$ is pinned to the wall at $A$ and held up by a tie
rod from a higher wall anchor $C$ to the boom tip $B$, where a load $W$
hangs. (b) The free-body diagram of the boom treated as one body. The
pin at $A$ supplies two reaction components $A_x$ and $A_y$; the tie rod,
a two-force member, pulls along its own axis with tension $T$ at angle
$\beta$ above the boom; the load $W$ hangs from the tip. Summing moments
about $A$ kills both pin components and solves for $T$ directly, after
which force balance returns $A_x$ and $A_y$. The governing connection is
whichever of the pin, the tie-rod thread, or the wall anchor first
reaches its capacity under the resolved forces.}
\label{fig:vol03ch01:jib-crane}
\end{figure}
